\noindent\textbf{Elección del procedimiento.} El campo es estacionario, por lo que $\partial_tu=0$, pero una partícula cambia de posición. Se aplica $a=u\,du/dx$. Para el tiempo se separa $dt=dx/u(x)$: usar longitud dividida por velocidad final sería incorrecto.

\begin{enumerate}

\item Sustituir $x=L$ en $u=U_0(1+x/L)$.

\[\fitmath{\mathrm{uL}=2\cdot U0}\]

\[\fitmath{\mathrm{uL}=2\cdot (10)\approx 20\,\mathrm{m/s}}\]

\item Derivar el campo respecto a $x$ antes de evaluar la aceleración.

\[\fitmath{\mathrm{grad}=\frac{U0}{L}}\]

\[\fitmath{\mathrm{grad}=\frac{(10)}{(0.5)}\approx 20\,\mathrm{1/s}}\]

\item Multiplicar la velocidad local por su gradiente espacial.

\[\fitmath{\mathrm{aL}=\mathrm{uL}\cdot \mathrm{grad}}\]

\[\fitmath{\mathrm{aL}=(20)\cdot (20)\approx 400\,\mathrm{m/s^{2}}}\]

\item La integral es $(L/U_0)[\ln(1+x/L)]_0^L$.

\[\fitmath{t=\frac{L}{U0}\cdot \ln\left(2\right)}\]

\[\fitmath{t=\frac{(0.5)}{(10)}\cdot \ln\left(2\right)\approx 0.03465736\,\mathrm{s}}\]

\end{enumerate}
